\setcounter{section}{6}
\section{Funtes del derecho}

Nombre completo: Las fuentes del derecho administrativo. La jerarquía de las fuentes. La ley. Las Disposiciones del Gobierno con fuerza de ley: decreto-ley y decreto legislativo. El reglamento: concepto, clases y límites. Otras fuentes del derecho administrativo.

\localtableofcontents

\subsection{Funtes del derecho Administrativo} %3
    Se habla de las fuentes del derecho como el origen y procedencia del derecho y las normas jurídicas. El conocimiento de estas fuentes es necesario para tener derecho.

    \subsubsection{Fuentes del derecho administrativo y sus clases}
        La teoría general de fuentes se recoge en el Art.1 del Código civil.
        \begin{itemize}
            \item Las fuentes del ordenamiento jurídico son la ley, la costumbre y los principios generales del derecho
            \item CArecerán de validez las disposiciones que contradigan otra de rango superior
            \item La costumbre solo regirá en defecto de ley aplicable, siempre que no sea contraria a la moral o al orden público y que resulte probada
            \item Los principios generales del derecho se aplicarán en defecto de la ley o costumbre, sin perjuicio de su carácter informador del ordenamiento jurídico
            \item Las normas jurídicas contenidas en los tratados internacionales no serán de aplicación directa en España en tanto no hayan pasado a formar parte del ordenamiento interno mediante su publicación íntegra en el «Boletín Oficial del Estado»
            \item La jurisprudencia complementará el ordenamiento jurídico con la doctrina que, de modo reiterado, establezca el Tribunal Supremo al interpretar y aplicar la ley, la costumbre y los principios generales del derecho
            \item Los Jueces y Tribunales tienen el deber inexcusable de resolver en todo caso los asuntos de que conozcan, ateniéndose al sistema de fuentes establecido.
        \end{itemize}
        Las fuenters pueden ser formales, cuando se refieren a la Ley, la costumbre y los Principios Generales del Derecho. Son Materiales cuando hay una institución con potestad normativa, como el Estado, las comunidades autónomas (leyes y reglamentos) o el pueblo (costumbre)
        %Según esto, la ley y la costumbre son leyes materiales y formales a la vez

        \paragraph{Clases}
            Según su proclamación son fuentes escritas o no escritas. Según el origen, pueden tener este en los poderes estatales, o en la costumbre. Es fuente primaria si se aplica en primer lugar y secundaria al resto. Las fuentes directas contienen normas y las indirectas no las tienen.

            Las fuentes indirectas no tienen normas jurídicas, pero ayudan a comprender y exteriorizar. Son ejemplos los tratados internacionales, y las doctrinas científicas o de los autores.

            El \textbf{Precedente administrativo} se refiere a que supuestos similares resueltos con la misma norma, deberían resultar igual. Aunque no hay una vinculación entre actuaciones previas, debe ser tenido en cuenta.

\subsection{Jerarquía de Fuentes} %6
    El orden jerárquico determina el orden de aplicación de las normas, y el criterio para solucionar las prescripciones contradictorias en normas de distinto rango. 
    
    \subsubsection{Principio de Jerarquía normativa}
        La Consitución lo reconoce en el Art 9.3:
        \newline
        \noindent\fbox{\parbox{\textwidth}{La Constitución garantiza el principio de legalidad, la jerarquía normativa, la publicidad de las normas, la irretroactividad de las disposiciones sancionadoras no favorables o restrictivas de derechos individuales, la seguridad jurídica, la responsabilidad y la interdicción de la arbitrariedad de los poderes públicos}}

        También en el código civil en su Art 1.2:
        \newline
        \noindent\fbox{\parbox{\textwidth}{
            Carecerán de validez las disposiciones que contradigan otra de rango superior.
        }}
        Igualmente en la ley 39/2015 de procedimiento administrativo en su Art 47.2:
        \newline
        \noindent\fbox{\parbox{\textwidth}{
            También serán nulas de pleno derecho las disposiciones administrativas que viulneren la constitución, las leyes u otras disposiciones administrativas de rango superior, las que regulen materias reservadas a la Ley, y las que establezcan la retroactividad de disposiciones sancionadoras no favorables o restrictivas de derechos individuales
        }}

        Es decir, las normas con rango superior prevalecen sobre las de rango inferior. Y en igual rango tienen la misma fuerza. Corresponde a la Constitución la norma suprema que prevalece sobre las demás, y que determina las potestades normativas. Todas las normas que provengan de la misma fuente tienen el mismo rango.

\subsection{La Ley} %8
    Ley es toda norma jurídica dictada por los órganos a los que el ordenamiento jurídico le atribuye poder legislativo. Según la constitución, las normas se han enriquecido, existiendo leyes orgánicas, ordinarias, de las CCAA, leyes marco, de transferencia, de armonización y normas con fuerza de ley como los Decretos Ley y Decrtos Legislativos.

    En un sentido ámplio se asimila a una norma jurídica o Derecho. Podría restringirse a normas jurídicas escritas. Y aún mas a norma de origen parlamentario, de rango subordinado a la Constitución e irresistible e indiscutible en el ámbito jurídico.

    Administrativamente, la ley es la norma jurídica de caracter general y obligatorio emanada del órgano estatal titular del Poder Legislativo. Estableciendo la parte formal (norma jurídica) y material (emanada del legislativo)

    \subsubsection{Características de la Ley}
        Alguinas características del término ley dentro del derecho son:
        \begin{description}
            \item[Obligatoriedad] Para todas las personas de respetar y cumplir, incluso contra su voluntad. El Art 6.1 del CC «La ignorancia de la ley no excusa de su cumplimiento»
            \item[Impersonabilidad] Se aplica a un número indeterminado de personas
            \item[Abstracción] Se aplican a todos los casos con los mismos supuestos
            \item[Permanencia] Se formulan con carácter indefinido y permanente. Hasta la subrogación, abrogación o derogación
            \item[Generalidad] Se aplican a todos los individuos cpon las condiciones determinadas
            \item[Irretroactividad] la retroactividad puede tener lugar cuando al ley lo disponga (retractividad expresa) o cuando se entienda en el sentido de la norma (retroactividad tácita). El TS declara que en el Art 2.3 del CC solo es válida la retroactividad expresa. Y la Constitución limita los casos de retroactividad.
        \end{description}
        Las leyes también tienen que ser imperativas (contener un mandato), Coercibles (que el mandato sea impuesto por coacción externa) y Racionales.

        Externamente, deben ser elaboradas y aprobadas según procedimiento establecido, además de sancionadas, promulgadas y publicadas.
    
    \subsubsection{Clases de leyes}
        Existen dos grandes fgrupos de leyes, las leyes orgánicas y las \st{demás} leyes ordinarias. Las leyes orgánicas tienen un contenido material limitado a las materias que la constitución exige: desarrollo de derechos fundamentales y libertades públicas, estatutos de autonomía, régimen electoral general y otras previstas, generalmente artículos cuyo desarrollo ulterior se deja a una ley orgánica.

        Las leyes orgánicas necesitan mayoría absoluta, contra la mayoría símple de las leyes ordinarias.

        La constitución establece otros tipos de leyes con alguna especialidad en su elaboración o aplicación:
        \begin{description}
            \item[Ley de Bases]delegación de las cortes al gobierno.
            \item[Leyes Marco] atribución a las comunidades autónomas para dictar en materias de competencia estatal para si mismas.
            \item[Leyes de transferencia] Leyes orgánicas en las que el estado transfiere facultades a las Comunidades Autónomas de las cuales el titular es el estado
            \item[Leyes de Armonización] establecen los principios para armonizar disposiciones normativas de las Comunidades Autónomas
            \item[Legislación Básica] leyes que contienen los principios y aspectos fundamentales de la regulación de una materia.
            \item[Ley de presupuestos] La elaboración corresponde exclusivamente al gobierno. Con un desarrollo singular, no pueden crear tributos, ni regular materias que no tengan relación con ingresos y gastos
        \end{description}

    \subsubsection{Leyes Orgánicas}
        Según el Art. 81 de la CE:
        \newline\noindent\fbox{\parbox{\textwidth}{
            Son leyes orgánicas las relativas al desarrollo de los derechos fundamentales y de las libertades públicas, las que aprueban los estatutos de autonomía y el régimen electoral general y las demás previstas en la constitución.
        }}

        Tienen una relación de superioridad sobre la ley ordinaria por dos principios. Por el de jerarquía normativa, se basa en la rigidez formal del Art.81.2. Por el principio de competencia por razón de materia, se acotan y separan ámbitos competenciales diferentes cuyo tratamiento se reserva a órganos y procedimientos determinados.

        En cuanto a los estatutos de Autonomía, el Art.147.1 establece que el Estado los reconoce y amparará como parte integrante de su Ordenamiento Jurídico. Los estatutos son jerárquicamente inferiores a la constitución. 

    \subsubsection{Leyes Ordinarias}
        Se aprueban por mayoría simple de los miembros de las Cámaras. Normalmente por el pleno, aunque se puede delegar en Comisiones Legislativas Permanentes. Quedan excluidas de la delegación la reforma constitucional, las cuestiones internacionales, las leyes orgánicas y de Bases y los Presupuestos generales.

        El procedimiento de elaboración puede iniciarse por las cámaras (proposición de ley) o por el Consejo de Ministros (proyecto de ley). La labor del senado en las leyes es meramente de revisión, ya que el congreso puede no atender a las propuestas.

\subsection{Disposiciones del Ejecutivo con fuerza de Ley} %12
    El gobierno dicta normas de rango reglamentario como parte de su ámbito reglamentativo propio. También aprueba normas con rango de ley, a pesar de no ser el representante de la voluntad popular.

    Las dos normas son los \textbf{Reales Decretos Legislativos} y \textbf{Reales Decretos Ley}. Los primeros se dictar por el gobierno para legislar una materia previa delegación del Poder Legislativo (El congreso le pide al gobierno que legisle) y los Reales Decretos Legislativos se dictan para afrontar situaciones de urgente necesidad. Ambas normas emanan del gobierno y tienen el mismo valor que las leyes ordinarias.

    La diferencia está en que el Real Decreto Legislativo, las cortes intervienen a priori, legitimando su promulgación y en los Reales Decretos Ley se interviene después.

    \subsubsection{Real Decreto Legislativo}
        El Art.82 permite a las Cortes Generales, delegar la postestad de dictar normas con rango de ley sobre materias no reservadas a leyes orgánicas.

        Consiste en una transferencia de la potestad legislativa en favor de la Administración, considerando una transferencia del ejercicio y no de la titularidad. El efecto primario es ampliar el poder normativo de la Administración. Esto se puede manifestar mediante \textbf{Textos Articulados de Leyes de Bases} o \textbf{Textos Refundidos}.

        \begin{itemize}
            \item Se otorga el poder en cada caso concreta y con un plazo para su ejercicio
            \item La delegación termina al publicarse la norma correspondiente
            \item No se concede de modo implícito ni por tiempo indeterminado
            \item No permite la subdelegación o a autoridades distintas del gobierno
            \item El destintario puede ser un gobierno autonómico si lo delega la asamblea autonómica.
            \item Esta delegación se hace mediante una ley de bases en los textos articulados, o una ley ordinaria cuando sea un texto refundido
        \end{itemize}
        Los Reales Decretos Legislativos pueden recibir tres tipos de control: los que establezcan las leyes de delegación expresamente, el del tribunal Constitucional y (según doctrina constitucional) las partes del real decreto que se excedan de los límites de la delegación, pierden el rango legal, quedando como reglamento y pudiendo ser revisada por la jurisdicción contencioso-administrativa.

    \subsubsection{Real Decreto Ley}
        Son normas emanadas del gobierno por motivos de extraordinaria y urgente necesidad. Estas normas son leyes desde el momento en el que se dictan, y se aprueban a posteriori en un plazo de 30 días. Esta convalidación convierte un Real Decreto Ley provisional en un Real Decreto Ley definitivo.

        El concepto de urgencia es subjetivo y discrecional, jurídicamente es indeterminado. El control de esta necesidad es político y no jurídico.

\subsection{El Reglamento: concepto, clases y límites} %15
    El reglamento es una norma escrita, pero es inferior a la ley. Su valor es subrogado a esta. No puede derogar una ley, y cualquier ley puede derogar un reglamento. No hay materias reservadas a los reglamentos. La ley puede entrar a regular cualquiera que haya sido regulada anteriormente en un reglamento.

    Son normas de caracter general, proceden del ejecutivo. Se dictan en el ejercicio de una competencia de los órganos administrativos, no son legislación delegada. Son normas de derecho objetivo que se integran en la legalidad y se imponen a la Administración en su actuación concreta. Pueden adoptar múltiple sformas como los Decretos, Órdenes Ministeriales\dots 

    Hay reserva material en los supuestos que la constitución exige que se regulen por ley, y formal, si el objeto ya está regulado por una ley.

    Los reglamentos pueden ser:
    \begin{description}
        \item[Independientes] Son dictados sin una ley previa
        \item[Ejecutivos] Dictados sobre una ley preexistente en virtud de la técnica de revisión normativa
        \item[De Necesidad] Dictado por la administración para hacer frente a situaciones extraordinarias.
    \end{description}
\subsection{Otras fuentes del derecho administrativo} %23
    \subsubsection{La Costumbre}
        Es la expresión del sentir popular o modelo de conducta popular expresado por la repeticióön de una serie de actos sobre un orden determinado de derecho. Es una norma creada por el uso social. Según el Código civil es necesario que no sea contraria a la moral o al orden público y que sea probada. Hay por tanto, un elemento material y un elemento espiritual.

        Hay diferentes tipos de costumbres:
        \begin{description}
            \item[\textit{Secundum legem}] regula una materia que ha sido objeto de ley, amplianda o dándole una interpretación peculiar
            \item[\textit{Praeter legem}] regula una materia que no ha sido regulada por ley
            \item[\textit{contra legem}] regula una materia distinta a como aparece regulada en la ley
        \end{description}

        Es un derecho fundamentalmente positivista y tiene un valor limitado.

    \subsubsection{Principios Generales del Derecho}
        se aplican en defecto de ley o costumbre, sin perjuicio de su caracter informador del ordenamiento jurídico. Estos se expresan «los valores materiales básicos de un ordenamiento jurídico, aquellos sobre los cuales se constiture como tal, las convicciones ético jurídicas fundamentales de una comunidad». No son ideas vagas o tendencias, sino principios del derecho.

    \subsubsection{Jurisprudencia}
        El cometido de la jurisprudencia es complementar al ordenamiento jurídico con la doctrina que se establezca al interpretar y aplicar la ley, la costumbre, y los principios generales del derecho por el Tribunal Supremo. En la Unión Europea estos son el \textbf{Tribunal Europeo de Derechos Humanos} y el \textbf{Tribunal de Justicia de la Comunidad Europea}.
% 26